% --- Archon physics formalization source begin ---
% archon:physics
% archon:covers PhyXMiniProblems/problem_phyx_mini_0123.lean
% archon:source-report reports/phyx_mini/problem_phyx_mini_0123.source.json

\chapter{Physics problem phyx\_mini\_0123}
\label{ch:PhyXMiniProblems_problem_phyx_mini_0123}

\paragraph{Problem source.}
\#\# Physical scenario

Quasars, an abbreviation for quasi-stellar radio sources, are distant objects that look like stars through a telescope but that emit far more electromagnetic radiation than an entire normal galaxy of stars. An example is the bright object below and to the left of center in figure; the other elongated objects in this image are normal galaxies. The leading model for the structure of a quasar is a galaxy with a supermassive black hole at its center. In this model, the radiation is emitted by interstellar gas and dust within the galaxy as this material falls toward the black hole. The radiation is thought to emanate from a region just a few light-years in diameter. (The diffuse glow surrounding the bright quasar shown in figure is thought to be this quasar's host galaxy.) To investigate this model of quasars and to study other exotic astronomical objects, the Russian Space Agency has placed a radio telescope in a large orbit around the earth. When this telescope is 77,000 km from earth and the signals it receives are combined with signals from the ground-based telescopes of the VLBA, the resolution is that of a single radio telescope 77,000 km in diameter. This arrangement can resolve in quasar 3C 405, which is \$7.2 \textbackslash{}times 10\textasciicircum{}8\$ light-years from earth, using radio waves at a frequency of 1665 MHz. Give your answer in light-years and in kilometers.

\#\# Question

What is the size of the smallest detail?

\#\# Answer choices

- A: A: \$1.94 \textbackslash{}times 10\textasciicircum{}\{13\}\$ km

- B: B: \$1.93 \textbackslash{}times 10\textasciicircum{}\{12\}\$ km

- C: C: \$1.95 \textbackslash{}times 10\textasciicircum{}\{14\}\$ km

- D: D: \$1.92 \textbackslash{}times 10\textasciicircum{}\{11\}\$ km

\#\# Figure caption (auxiliary; use the image as primary evidence)

This image appears to be a grayscale photograph of a celestial object, likely taken through a telescope. Key features include:

- **Central Bright Object**: Dominating the image is a bright point of light at the center, which could be a star or the core of a galaxy. It has a noticeable glare or halo around it.
  
- **Peripheral Structures**: Around the central object, there are faint, elongated structures, possibly other galaxies or cosmic features, extending in various directions. They appear to be distributed unevenly across the image.

- **Noise/Grain**: The image has a significant amount of noise or grain, typical of deep-space photographs taken with high sensitivity to capture faint details.

- **No Text**: There is no visible text present within the image itself.

The scene captures a section of space with multiple celestial objects, highlighting the contrast between a bright central source and fainter surrounding structures.

\#\# Dataset metadata

Optical Instruments; Physical Model Grounding Reasoning, Numerical Reasoning

\paragraph{Recorded answer/context.}
A: A: \$1.94 \textbackslash{}times 10\textasciicircum{}\{13\}\$ km

\paragraph{Figure/image path.}
/root/proposal\_for\_physic/science-mango/phyx\_mini\_run/phyx\_data/test\_image/123.png

\paragraph{Formalization target.}
create a compiling Lean file with sorry bodies at `PhyXMiniProblems/problem\_phyx\_mini\_0123.lean`. The Lean declarations must preserve the physical quantities, dimensions or dimensional roles, figure labels, governing-law hypotheses, and final relation expressed by this problem.
Use Mathlib/Physlib names found through LeanExplore where available. If a physics API is missing, introduce faithful local abstractions rather than scalar placeholder aliases.

\begin{theorem}[Physics formalization target]
\label{thm:physics:phyx_mini_0123:target}
\lean{PhyXMiniProblems.ProblemPhyXMini0123.problem_phyx_mini_0123}
\uses{def:physics:phyx-mini-0123:phyxminiproblems-problemphyxmini0123-quasarresolutionsetup, def:physics:phyx-mini-0123:phyxminiproblems-problemphyxmini0123-matchesscenarioandfigure, def:physics:phyx-mini-0123:phyxminiproblems-problemphyxmini0123-matchesproblemreadouts, def:physics:phyx-mini-0123:phyxminiproblems-problemphyxmini0123-hasphysicalparameters, def:physics:phyx-mini-0123:phyxminiproblems-problemphyxmini0123-satisfiesinterferometricaperturemodel, def:physics:phyx-mini-0123:phyxminiproblems-problemphyxmini0123-satisfiesvacuumwavelaw, def:physics:phyx-mini-0123:phyxminiproblems-problemphyxmini0123-satisfiesrayleighcriterion, def:physics:phyx-mini-0123:phyxminiproblems-problemphyxmini0123-satisfiesparaxialresolutiongeometry, def:physics:phyx-mini-0123:phyxminiproblems-problemphyxmini0123-matchesdisplayedkilometerchoice, def:physics:phyx-mini-0123:phyxminiproblems-problemphyxmini0123-roundstolightyearanswer}
The assigned autoformalize agent should translate this physics problem into Lean declarations in the covered file, with theorem and lemma proof bodies written as `by sorry`.
\end{theorem}
\begin{proof}
This is an autoformalization task, not a proof task. Produce statements that can later be proved by the physics prover without weakening the physical model.
\end{proof}

% --- Archon declaration topology begin ---
\section{Declaration topology}
\begin{definition}[Length Quantity]
  \label{def:physics:phyx-mini-0123:phyxminiproblems-problemphyxmini0123-lengthquantity}
  \lean{PhyXMiniProblems.ProblemPhyXMini0123.LengthQuantity}
  A nonnegative physical length, independent of the unit used to read it
\end{definition}
\begin{proof}
  This is the declared model definition of Length Quantity.
\end{proof}
\begin{definition}[Frequency Quantity]
  \label{def:physics:phyx-mini-0123:phyxminiproblems-problemphyxmini0123-frequencyquantity}
  \lean{PhyXMiniProblems.ProblemPhyXMini0123.FrequencyQuantity}
  A nonnegative physical frequency, carrying inverse-time dimension
\end{definition}
\begin{proof}
  This is the declared model definition of Frequency Quantity.
\end{proof}
\begin{definition}[length Readout]
  \label{def:physics:phyx-mini-0123:phyxminiproblems-problemphyxmini0123-lengthreadout}
  \lean{PhyXMiniProblems.ProblemPhyXMini0123.lengthReadout}
  \uses{def:physics:phyx-mini-0123:phyxminiproblems-problemphyxmini0123-lengthquantity}
  Scalar readout of a physical length in a selected length unit
\end{definition}
\begin{proof}
  This is the declared model definition of length Readout.
\end{proof}
\begin{definition}[frequency Readout]
  \label{def:physics:phyx-mini-0123:phyxminiproblems-problemphyxmini0123-frequencyreadout}
  \lean{PhyXMiniProblems.ProblemPhyXMini0123.frequencyReadout}
  \uses{def:physics:phyx-mini-0123:phyxminiproblems-problemphyxmini0123-frequencyquantity}
  Scalar readout of a physical frequency in inverse selected time units
\end{definition}
\begin{proof}
  This is the declared model definition of frequency Readout.
\end{proof}
\begin{definition}[length In Meters]
  \label{def:physics:phyx-mini-0123:phyxminiproblems-problemphyxmini0123-lengthinmeters}
  \lean{PhyXMiniProblems.ProblemPhyXMini0123.lengthInMeters}
  \uses{def:physics:phyx-mini-0123:phyxminiproblems-problemphyxmini0123-lengthquantity, def:physics:phyx-mini-0123:phyxminiproblems-problemphyxmini0123-lengthreadout}
  Meter readout used for positivity and the SI wave relation
\end{definition}
\begin{proof}
  This is the declared model definition of length In Meters.
\end{proof}
\begin{definition}[length In Kilometers]
  \label{def:physics:phyx-mini-0123:phyxminiproblems-problemphyxmini0123-lengthinkilometers}
  \lean{PhyXMiniProblems.ProblemPhyXMini0123.lengthInKilometers}
  \uses{def:physics:phyx-mini-0123:phyxminiproblems-problemphyxmini0123-lengthquantity, def:physics:phyx-mini-0123:phyxminiproblems-problemphyxmini0123-lengthreadout}
  Kilometer readout used by the telescope baseline and answer choices
\end{definition}
\begin{proof}
  This is the declared model definition of length In Kilometers.
\end{proof}
\begin{definition}[length In Light Years]
  \label{def:physics:phyx-mini-0123:phyxminiproblems-problemphyxmini0123-lengthinlightyears}
  \lean{PhyXMiniProblems.ProblemPhyXMini0123.lengthInLightYears}
  \uses{def:physics:phyx-mini-0123:phyxminiproblems-problemphyxmini0123-lengthquantity, def:physics:phyx-mini-0123:phyxminiproblems-problemphyxmini0123-lengthreadout}
  Light-year readout used for the source distance and one requested answer
\end{definition}
\begin{proof}
  This is the declared model definition of length In Light Years.
\end{proof}
\begin{definition}[frequency In Hertz]
  \label{def:physics:phyx-mini-0123:phyxminiproblems-problemphyxmini0123-frequencyinhertz}
  \lean{PhyXMiniProblems.ProblemPhyXMini0123.frequencyInHertz}
  \uses{def:physics:phyx-mini-0123:phyxminiproblems-problemphyxmini0123-frequencyquantity, def:physics:phyx-mini-0123:phyxminiproblems-problemphyxmini0123-frequencyreadout}
  Hertz readout of a physical frequency
\end{definition}
\begin{proof}
  This is the declared model definition of frequency In Hertz.
\end{proof}
\begin{definition}[vacuum Speed Of Light Readout]
  \label{def:physics:phyx-mini-0123:phyxminiproblems-problemphyxmini0123-vacuumspeedoflightreadout}
  \lean{PhyXMiniProblems.ProblemPhyXMini0123.vacuumSpeedOfLightReadout}
  Physlib's exact vacuum speed of light in selected length and time units
\end{definition}
\begin{proof}
  This is the declared model definition of vacuum Speed Of Light Readout.
\end{proof}
\begin{definition}[Astronomical Source Label]
  \label{def:physics:phyx-mini-0123:phyxminiproblems-problemphyxmini0123-astronomicalsourcelabel}
  \lean{PhyXMiniProblems.ProblemPhyXMini0123.AstronomicalSourceLabel}
  The astronomical source named in the problem
\end{definition}
\begin{proof}
  This is the declared model definition of Astronomical Source Label.
\end{proof}
\begin{definition}[Astronomical Source Class]
  \label{def:physics:phyx-mini-0123:phyxminiproblems-problemphyxmini0123-astronomicalsourceclass}
  \lean{PhyXMiniProblems.ProblemPhyXMini0123.AstronomicalSourceClass}
  Physical class assigned to the source in the scenario
\end{definition}
\begin{proof}
  This is the declared model definition of Astronomical Source Class.
\end{proof}
\begin{definition}[Quasar Emission Model]
  \label{def:physics:phyx-mini-0123:phyxminiproblems-problemphyxmini0123-quasaremissionmodel}
  \lean{PhyXMiniProblems.ProblemPhyXMini0123.QuasarEmissionModel}
  Leading qualitative model for the quasar's radiation source
\end{definition}
\begin{proof}
  This is the declared model definition of Quasar Emission Model.
\end{proof}
\begin{definition}[Emission Region Scale]
  \label{def:physics:phyx-mini-0123:phyxminiproblems-problemphyxmini0123-emissionregionscale}
  \lean{PhyXMiniProblems.ProblemPhyXMini0123.EmissionRegionScale}
  Qualitative scale stated for the compact emitting region
\end{definition}
\begin{proof}
  This is the declared model definition of Emission Region Scale.
\end{proof}
\begin{definition}[Radio Array Configuration]
  \label{def:physics:phyx-mini-0123:phyxminiproblems-problemphyxmini0123-radioarrayconfiguration}
  \lean{PhyXMiniProblems.ProblemPhyXMini0123.RadioArrayConfiguration}
  Radio-instrument configuration used for the observation
\end{definition}
\begin{proof}
  This is the declared model definition of Radio Array Configuration.
\end{proof}
\begin{definition}[Resolution Regime]
  \label{def:physics:phyx-mini-0123:phyxminiproblems-problemphyxmini0123-resolutionregime}
  \lean{PhyXMiniProblems.ProblemPhyXMini0123.ResolutionRegime}
  Optical regime used to estimate the smallest resolvable detail
\end{definition}
\begin{proof}
  This is the declared model definition of Resolution Regime.
\end{proof}
\begin{definition}[Figure Feature]
  \label{def:physics:phyx-mini-0123:phyxminiproblems-problemphyxmini0123-figurefeature}
  \lean{PhyXMiniProblems.ProblemPhyXMini0123.FigureFeature}
  Distinct qualitative objects visible in the supplied astronomical image
\end{definition}
\begin{proof}
  This is the declared model definition of Figure Feature.
\end{proof}
\begin{definition}[Image Location]
  \label{def:physics:phyx-mini-0123:phyxminiproblems-problemphyxmini0123-imagelocation}
  \lean{PhyXMiniProblems.ProblemPhyXMini0123.ImageLocation}
  Qualitative locations distinguished by the astronomical photograph
\end{definition}
\begin{proof}
  This is the declared model definition of Image Location.
\end{proof}
\begin{definition}[Quasar Resolution Setup]
  \label{def:physics:phyx-mini-0123:phyxminiproblems-problemphyxmini0123-quasarresolutionsetup}
  \lean{PhyXMiniProblems.ProblemPhyXMini0123.QuasarResolutionSetup}
  \uses{def:physics:phyx-mini-0123:phyxminiproblems-problemphyxmini0123-lengthquantity, def:physics:phyx-mini-0123:phyxminiproblems-problemphyxmini0123-frequencyquantity, def:physics:phyx-mini-0123:phyxminiproblems-problemphyxmini0123-astronomicalsourcelabel, def:physics:phyx-mini-0123:phyxminiproblems-problemphyxmini0123-astronomicalsourceclass, def:physics:phyx-mini-0123:phyxminiproblems-problemphyxmini0123-quasaremissionmodel, def:physics:phyx-mini-0123:phyxminiproblems-problemphyxmini0123-emissionregionscale, def:physics:phyx-mini-0123:phyxminiproblems-problemphyxmini0123-radioarrayconfiguration, def:physics:phyx-mini-0123:phyxminiproblems-problemphyxmini0123-resolutionregime, def:physics:phyx-mini-0123:phyxminiproblems-problemphyxmini0123-figurefeature, def:physics:phyx-mini-0123:phyxminiproblems-problemphyxmini0123-imagelocation}
  All physical quantities and qualitative labels in the problem. smallestResolvableDetail is an unknown physical length. No field gives it a numerical value or identifies an answer choice
\end{definition}
\begin{proof}
  This is the declared model definition of Quasar Resolution Setup.
\end{proof}
\begin{definition}[Matches Scenario And Figure]
  \label{def:physics:phyx-mini-0123:phyxminiproblems-problemphyxmini0123-matchesscenarioandfigure}
  \lean{PhyXMiniProblems.ProblemPhyXMini0123.MatchesScenarioAndFigure}
  \uses{def:physics:phyx-mini-0123:phyxminiproblems-problemphyxmini0123-quasarresolutionsetup}
  Qualitative scenario and primary-image readouts. The image contains the bright compact quasar, diffuse host-galaxy glow, and elongated normal galaxies; it supplies no numerical length scale
\end{definition}
\begin{proof}
  This is the declared model definition of Matches Scenario And Figure.
\end{proof}
\begin{definition}[Matches Problem Readouts]
  \label{def:physics:phyx-mini-0123:phyxminiproblems-problemphyxmini0123-matchesproblemreadouts}
  \lean{PhyXMiniProblems.ProblemPhyXMini0123.MatchesProblemReadouts}
  \uses{def:physics:phyx-mini-0123:phyxminiproblems-problemphyxmini0123-lengthinkilometers, def:physics:phyx-mini-0123:phyxminiproblems-problemphyxmini0123-lengthinlightyears, def:physics:phyx-mini-0123:phyxminiproblems-problemphyxmini0123-frequencyinhertz, def:physics:phyx-mini-0123:phyxminiproblems-problemphyxmini0123-quasarresolutionsetup}
  Numerical data explicitly stated in the problem: a 77,000 km space-to-Earth baseline, distance 7.2 * 10\textasciicircum{}8 light-years, and radio frequency 1665 MHz. The latter is recorded as 1665 * 10\textasciicircum{}6 Hz
\end{definition}
\begin{proof}
  This is the declared model definition of Matches Problem Readouts.
\end{proof}
\begin{definition}[Has Physical Parameters]
  \label{def:physics:phyx-mini-0123:phyxminiproblems-problemphyxmini0123-hasphysicalparameters}
  \lean{PhyXMiniProblems.ProblemPhyXMini0123.HasPhysicalParameters}
  \uses{def:physics:phyx-mini-0123:phyxminiproblems-problemphyxmini0123-lengthinmeters, def:physics:phyx-mini-0123:phyxminiproblems-problemphyxmini0123-frequencyinhertz, def:physics:phyx-mini-0123:phyxminiproblems-problemphyxmini0123-quasarresolutionsetup}
  Positivity and angular-range conditions for the physical configuration
\end{definition}
\begin{proof}
  This is the declared model definition of Has Physical Parameters.
\end{proof}
\begin{definition}[Satisfies Interferometric Aperture Model]
  \label{def:physics:phyx-mini-0123:phyxminiproblems-problemphyxmini0123-satisfiesinterferometricaperturemodel}
  \lean{PhyXMiniProblems.ProblemPhyXMini0123.SatisfiesInterferometricApertureModel}
  \uses{def:physics:phyx-mini-0123:phyxminiproblems-problemphyxmini0123-lengthreadout, def:physics:phyx-mini-0123:phyxminiproblems-problemphyxmini0123-quasarresolutionsetup}
  The problem's interferometric idealization: combining the orbiting receiver with the VLBA gives the resolution of one circular aperture whose diameter is the space-to-Earth baseline. This is required in every length unit
\end{definition}
\begin{proof}
  This is the declared model definition of Satisfies Interferometric Aperture Model.
\end{proof}
\begin{definition}[Satisfies Vacuum Wave Law]
  \label{def:physics:phyx-mini-0123:phyxminiproblems-problemphyxmini0123-satisfiesvacuumwavelaw}
  \lean{PhyXMiniProblems.ProblemPhyXMini0123.SatisfiesVacuumWaveLaw}
  \uses{def:physics:phyx-mini-0123:phyxminiproblems-problemphyxmini0123-lengthreadout, def:physics:phyx-mini-0123:phyxminiproblems-problemphyxmini0123-frequencyreadout, def:physics:phyx-mini-0123:phyxminiproblems-problemphyxmini0123-vacuumspeedoflightreadout, def:physics:phyx-mini-0123:phyxminiproblems-problemphyxmini0123-quasarresolutionsetup}
  Vacuum electromagnetic-wave dispersion, c = λ f, in every compatible choice of length and time units. This determines the wavelength but not the requested resolvable size
\end{definition}
\begin{proof}
  This is the declared model definition of Satisfies Vacuum Wave Law.
\end{proof}
\begin{definition}[Satisfies Rayleigh Criterion]
  \label{def:physics:phyx-mini-0123:phyxminiproblems-problemphyxmini0123-satisfiesrayleighcriterion}
  \lean{PhyXMiniProblems.ProblemPhyXMini0123.SatisfiesRayleighCriterion}
  \uses{def:physics:phyx-mini-0123:phyxminiproblems-problemphyxmini0123-lengthreadout, def:physics:phyx-mini-0123:phyxminiproblems-problemphyxmini0123-quasarresolutionsetup}
  Rayleigh's circular-aperture criterion θ = 1.22 λ / D, with the radian angle dimensionless. The factor 1.22 is represented exactly by 61 / 50
\end{definition}
\begin{proof}
  This is the declared model definition of Satisfies Rayleigh Criterion.
\end{proof}
\begin{definition}[Satisfies Paraxial Resolution Geometry]
  \label{def:physics:phyx-mini-0123:phyxminiproblems-problemphyxmini0123-satisfiesparaxialresolutiongeometry}
  \lean{PhyXMiniProblems.ProblemPhyXMini0123.SatisfiesParaxialResolutionGeometry}
  \uses{def:physics:phyx-mini-0123:phyxminiproblems-problemphyxmini0123-lengthreadout, def:physics:phyx-mini-0123:phyxminiproblems-problemphyxmini0123-quasarresolutionsetup}
  Small-angle geometry at the source: the physical linear detail is the source distance times the angular resolution. This is a governing relation, not a numerical assignment of the unknown detail
\end{definition}
\begin{proof}
  This is the declared model definition of Satisfies Paraxial Resolution Geometry.
\end{proof}
\begin{definition}[Answer Choice]
  \label{def:physics:phyx-mini-0123:phyxminiproblems-problemphyxmini0123-answerchoice}
  \lean{PhyXMiniProblems.ProblemPhyXMini0123.AnswerChoice}
  Labels of the four kilometer-valued choices printed with the problem
\end{definition}
\begin{proof}
  This is the declared model definition of Answer Choice.
\end{proof}
\begin{definition}[displayed Detail Kilometers]
  \label{def:physics:phyx-mini-0123:phyxminiproblems-problemphyxmini0123-displayeddetailkilometers}
  \lean{PhyXMiniProblems.ProblemPhyXMini0123.displayedDetailKilometers}
  \uses{def:physics:phyx-mini-0123:phyxminiproblems-problemphyxmini0123-answerchoice}
  Kilometer readout displayed beside each answer choice
\end{definition}
\begin{proof}
  This is the declared model definition of displayed Detail Kilometers.
\end{proof}
\begin{definition}[displayed Precision Kilometers]
  \label{def:physics:phyx-mini-0123:phyxminiproblems-problemphyxmini0123-displayedprecisionkilometers}
  \lean{PhyXMiniProblems.ProblemPhyXMini0123.displayedPrecisionKilometers}
  \uses{def:physics:phyx-mini-0123:phyxminiproblems-problemphyxmini0123-answerchoice}
  Place value of the last displayed significant digit, in kilometers
\end{definition}
\begin{proof}
  This is the declared model definition of displayed Precision Kilometers.
\end{proof}
\begin{definition}[recorded Dataset Answer]
  \label{def:physics:phyx-mini-0123:phyxminiproblems-problemphyxmini0123-recordeddatasetanswer}
  \lean{PhyXMiniProblems.ProblemPhyXMini0123.recordedDatasetAnswer}
  \uses{def:physics:phyx-mini-0123:phyxminiproblems-problemphyxmini0123-answerchoice}
  The answer label recorded in the source dataset
\end{definition}
\begin{proof}
  This is the declared model definition of recorded Dataset Answer.
\end{proof}
\begin{definition}[Matches Displayed Kilometer Choice]
  \label{def:physics:phyx-mini-0123:phyxminiproblems-problemphyxmini0123-matchesdisplayedkilometerchoice}
  \lean{PhyXMiniProblems.ProblemPhyXMini0123.MatchesDisplayedKilometerChoice}
  \uses{def:physics:phyx-mini-0123:phyxminiproblems-problemphyxmini0123-lengthinkilometers, def:physics:phyx-mini-0123:phyxminiproblems-problemphyxmini0123-quasarresolutionsetup, def:physics:phyx-mini-0123:phyxminiproblems-problemphyxmini0123-answerchoice, def:physics:phyx-mini-0123:phyxminiproblems-problemphyxmini0123-displayeddetailkilometers, def:physics:phyx-mini-0123:phyxminiproblems-problemphyxmini0123-displayedprecisionkilometers}
  Agreement of the derived detail with a displayed kilometer answer
\end{definition}
\begin{proof}
  This is the declared model definition of Matches Displayed Kilometer Choice.
\end{proof}
\begin{definition}[Rounds To Light Year Answer]
  \label{def:physics:phyx-mini-0123:phyxminiproblems-problemphyxmini0123-roundstolightyearanswer}
  \lean{PhyXMiniProblems.ProblemPhyXMini0123.RoundsToLightYearAnswer}
  \uses{def:physics:phyx-mini-0123:phyxminiproblems-problemphyxmini0123-lengthinlightyears, def:physics:phyx-mini-0123:phyxminiproblems-problemphyxmini0123-quasarresolutionsetup}
  The light-year answer rounds to 2.05 ly at the stated precision
\end{definition}
\begin{proof}
  This is the declared model definition of Rounds To Light Year Answer.
\end{proof}
% --- Archon declaration topology end ---
% --- Archon physics formalization source end ---
